\documentclass[11pt]{article}

% ACL 2025 style — use acl package from Overleaf or ACL Anthology
\usepackage{acl}

% Standard packages
\usepackage{times}
\usepackage{latexsym}
\usepackage{graphicx}
\usepackage{booktabs}
\usepackage{multirow}
\usepackage{amsmath}
\usepackage{amssymb}
\usepackage{url}
\usepackage{microtype}
\usepackage{xcolor}
\usepackage{tikz}
\usepackage{pgfplots}
\pgfplotsset{compat=1.18}

% Fix: define \cmark and \xmark for ablation table
\usepackage{pifont}
\newcommand{\cmark}{\ding{51}}% ✓
\newcommand{\xmark}{\ding{55}}% ✗

% -------------------------------------------------------
% Title and Authors
% -------------------------------------------------------
\title{TempoRAG-KG: Temporal Knowledge Graph-Augmented\\
Retrieval for Multi-Hop Question Answering}

\author{
  Supanut Kompayak \quad Aphisit Jaemyaem \quad Dechathon Niamsa-ard \quad Kaung Hein Htet \\
  Asian Institute of Technology, Pathum Thani, Thailand \\
  \texttt{\{st126055, st126130, st126235, st126477\}@ait.asia}
}

\begin{document}
\maketitle

% -------------------------------------------------------
% ABSTRACT
% -------------------------------------------------------
\begin{abstract}
Knowledge Graph-augmented Retrieval (KG-RAG) systems have demonstrated
strong performance on multi-hop question answering by structuring retrieved
facts as interconnected triples. However, existing systems---including the
recent state-of-the-art KG$^2$RAG---assume a static snapshot of knowledge
and fail to account for \textit{temporal validity}: the reality that facts
are only true within specific time windows. In this paper, we present
\textbf{TempoRAG-KG}, a temporally-aware KG-RAG framework that extends
KG$^2$RAG with two contributions: (1) explicit $[\mathtt{valid\_from},
\mathtt{valid\_to}]$ validity intervals attached to every KG triple,
enabling deterministic temporal filtering during retrieval, and (2) a
Graph-of-Graph (GoG)-based fill-in mechanism for handling incomplete KG
facts. We first conduct a systematic analysis of temporal failure modes in
KG-RAG, categorizing failures into four types and measuring their
distribution across HotpotQA and MuSiQue. Our exploratory data analysis
reveals that 35.0\% of HotpotQA and 38.3\% of MuSiQue development
questions contain temporal expressions, with prevalence increasing from
30.2\% in 2-hop to 46.7\% in 4-hop questions---suggesting temporal validity
becomes increasingly critical as reasoning complexity grows. Ablation
experiments show that TempoRAG-KG improves F1 over the KG$^2$RAG baseline,
with gains concentrated in temporally-sensitive question subsets.
\end{abstract}

% -------------------------------------------------------
% 1. INTRODUCTION
% -------------------------------------------------------
\section{Introduction}

% Paragraph 1 — State of the world
Retrieval-Augmented Generation (RAG) has emerged as the dominant paradigm
for grounding large language model (LLM) responses in external factual
knowledge \cite{lewis2020retrieval}. Recent advances have moved beyond
flat vector retrieval toward structured graph-based approaches, where
Knowledge Graphs (KGs) provide fact-level relationships between retrieved
chunks. The state-of-the-art KG$^2$RAG \cite{zhu2025kg2rag} demonstrates
that organizing retrieved evidence through KG-guided expansion and Maximum
Spanning Tree context organization yields substantial improvements over
prior RAG baselines on multi-hop question answering benchmarks such as
HotpotQA \cite{yang2018hotpotqa}.

% Paragraph 2 — The big BUT (most important)
However, a fundamental limitation remains unaddressed: \textbf{all existing
KG-RAG systems, including KG$^2$RAG, assume a static knowledge graph in
which every triple is treated as universally and permanently valid.} In
reality, facts are only true within specific time periods---Tim Cook
\textit{became} CEO of Apple in 2011, implying Steve Jobs held that role
before him. When a user asks \textit{``Who was the CEO of Apple in 2005?''},
a system without temporal awareness retrieves both facts simultaneously and
either hallucinates or selects arbitrarily. We term this class of errors
\textit{temporal failure modes}. Despite the prevalence of such queries---our
exploratory data analysis (EDA) reveals that 35.0\% of HotpotQA and 38.3\%
of MuSiQue development questions contain temporal expressions---no existing
KG-RAG framework systematically addresses temporal validity at the
triple level on general-domain unstructured text.

% Paragraph 3 — Therefore, we did
To address this gap, we propose \textbf{TempoRAG-KG}, a framework that
extends KG$^2$RAG with two targeted components. First, we attach explicit
$[\mathtt{valid\_from}, \mathtt{valid\_to}]$ validity intervals to every
extracted triple during offline KG construction, and apply deterministic
temporal filtering prior to graph expansion. Second, we integrate a
GoG-based fill-in mechanism \cite{he2024gog} to recover answers when the
constructed KG is incomplete. Crucially, we precede system development
with a \textit{systematic diagnostic study} that categorizes temporal
failure modes into four types and quantifies their distribution across
two established multi-hop QA benchmarks.

% Paragraph 4 — Key findings
Our analysis reveals several important findings. First, temporal failure
modes in KG-RAG are not uniformly distributed: stale-fact retrieval
accounts for the majority of temporal errors, while multi-hop temporal
chain failures are rarer but catastrophic for answer quality. Second,
temporal prevalence escalates with reasoning complexity---from 30.2\% in
2-hop MuSiQue questions to 46.7\% in 4-hop questions---suggesting that
temporal awareness is disproportionately important for complex multi-hop
chains. Third, LLM-based temporal extraction achieves high precision on
explicit date expressions but degrades substantially on implicit and
conflicting temporal references, identifying a clear bottleneck for
future work.

% Paragraph 5 — Contributions
The contributions of this work are:
\begin{enumerate}
  \item \textbf{Temporal Failure Taxonomy}: A systematic categorization of
        four temporal failure modes in KG-RAG systems, with empirical
        distribution analysis across HotpotQA and MuSiQue.
  \item \textbf{TempoRAG-KG Framework}: A temporally-aware KG-RAG pipeline
        that attaches explicit validity intervals per triple and applies
        temporal filtering prior to graph expansion, extending
        KG$^2$RAG \cite{zhu2025kg2rag}.
  \item \textbf{Temporal Extraction Study}: An empirical evaluation of
        LLM-based temporal validity extraction from unstructured text,
        characterizing accuracy across three linguistic pattern types.
\end{enumerate}

% -------------------------------------------------------
% 2. RELATED WORK
% -------------------------------------------------------
\section{Related Work}

\subsection{KG-Augmented Retrieval-Augmented Generation}

Early RAG systems \cite{lewis2020retrieval} rely on dense passage retrieval
to retrieve isolated text chunks, ignoring structural relationships between
facts. Graph-based RAG approaches address this by incorporating KG structure
into retrieval. KG$^2$RAG \cite{zhu2025kg2rag} proposes KG-guided chunk
expansion using seed retrieval followed by a KG-based chunk organization
process via Maximum Spanning Tree, achieving strong performance on
HotpotQA. GEAR \cite{ma2025gear} introduces efficient beam-search graph
traversal that replaces expensive LLM calls at each hop with cosine
similarity scoring, substantially reducing inference cost. K-RagRec
\cite{wang2025kragrec} applies KG-guided RAG to recommendation, achieving
up to 93\% hallucination reduction. Despite these advances,
\textbf{none of these systems attach temporal validity metadata to KG
triples}, treating all retrieved facts as equally and permanently valid.

\subsection{Temporal Knowledge Graph Reasoning}

Temporal KG reasoning has been studied extensively over pre-built structured
KGs such as Wikidata and ICEWS. TempAgent \cite{tempagent2025} introduces
temporal filtering tools for an LLM agent, but requires a pre-existing
structured temporal KG and cannot operate over general-domain unstructured
text. EvoReasoner \cite{evoreasonerevokg2025} proposes a unified framework
combining multi-hop temporal reasoning with incremental KG evolution from
documents via confidence-based contradiction resolution and temporal trend
tracking---the closest existing work to TempoRAG-KG. However, EvoReasoner
tracks temporal evolution \textit{implicitly} through confidence scores
rather than explicit $[\mathtt{valid\_from}, \mathtt{valid\_to}]$ intervals,
which prevents deterministic point-in-time querying. Furthermore,
EvoReasoner is evaluated on structured TKGQA benchmarks (MultiTQ,
CronQuestions) rather than general-domain multi-hop QA datasets.
\textbf{Our work is the first to combine explicit triple-level temporal
validity tagging with general-domain KG-RAG and evaluate on HotpotQA and
MuSiQue.}

\subsection{KG Construction from Unstructured Text}

Automatic KG construction enables RAG systems to operate without pre-built
KGs. HRGraph \cite{hrgraph2024} and STINMatch \cite{stinmatch2023} use LLM
prompting to extract subject-relation-object triples from raw text, while
TOBUGraph \cite{tobugraph2024} demonstrates production-scale construction.
GoG \cite{he2024gog} addresses KG incompleteness by allowing LLMs to fill
in missing facts from parametric knowledge when graph traversal reaches
dead ends. However, \textbf{none of these approaches extract or store
temporal validity intervals alongside triples}, losing the temporal context
necessary for time-sensitive queries. TempoRAG-KG extends LLM-based
extraction to additionally infer $\mathtt{valid\_from}$ and
$\mathtt{valid\_to}$ for each triple.

\paragraph{Gap Summary.}
No existing system simultaneously (1) automatically constructs a KG from
general-domain unstructured text, (2) attaches explicit per-triple temporal
validity intervals, and (3) evaluates on multi-hop QA benchmarks. TempoRAG-KG
fills this gap.

% -------------------------------------------------------
% Figure 2 placeholder — Temporal Failure Taxonomy
% -------------------------------------------------------
\begin{figure}[t]
  \centering
  \fbox{\parbox{0.45\textwidth}{
    \centering
    \vspace{1.5cm}
    \textbf{[Figure 2: Temporal Failure Taxonomy]}\\[0.5em]
    2$\times$2 grid of four failure types:\\[0.3em]
    \textbf{Type 1} Stale Fact: KG retrieves an expired fact\\
    Example: ``Steve Jobs'' retrieved as current Apple CEO\\[0.3em]
    \textbf{Type 2} Conflicting Facts: Multiple contradictory facts retrieved\\
    Example: Both Jobs and Cook returned as Apple CEO\\[0.3em]
    \textbf{Type 3} Missing Temporal Context: Correct entity, unknown time window\\
    Example: Cook is CEO but no valid\_from/valid\_to known\\[0.3em]
    \textbf{Type 4} Temporal Hop Failure: Any hop retrieves a stale fact\\
    Example: Hop 1 correct, hop 2 retrieves outdated fact
    \vspace{1.5cm}
  }}
  \caption{Four types of temporal failure modes in KG-RAG systems,
  identified through manual analysis of KG$^2$RAG errors on the
  HotpotQA temporal subset.}
  \label{fig:taxonomy}
\end{figure}

% -------------------------------------------------------
% 3. RESEARCH QUESTIONS
% -------------------------------------------------------
\section{Research Questions}

This work is organized around three research questions:

\begin{description}
  \item[\textbf{RQ1}] \textit{What types of temporal failures occur in
    KG-RAG systems when answering time-sensitive multi-hop questions,
    and how are they distributed across HotpotQA and MuSiQue?}\\
    \textbf{IV:} Question type (temporal vs.\ non-temporal), hop count.
    \textbf{DV:} Failure type distribution, failure rate per category.

  \item[\textbf{RQ2}] \textit{To what extent does explicit triple-level
    temporal validity tagging improve retrieval precision and answer F1
    on temporally-sensitive questions compared to the KG$^2$RAG baseline?}\\
    \textbf{IV:} Presence of temporal tagging and filtering.
    \textbf{DV:} F1 score, Exact Match (EM), retrieval precision on
    temporal vs.\ non-temporal question subsets.

  \item[\textbf{RQ3}] \textit{How accurately can LLMs extract temporal
    validity intervals from unstructured text, and what linguistic patterns
    most affect extraction quality?}\\
    \textbf{IV:} Temporal expression type (explicit year, implicit/relative,
    conflicting).
    \textbf{DV:} Extraction precision, extraction recall, F1.
\end{description}

% -------------------------------------------------------
% 4. METHODOLOGY
% -------------------------------------------------------
\section{Methodology}

% -------------------------------------------------------
% Figure 1 — System Architecture (placed here so it floats to top of Methodology page)
% -------------------------------------------------------
\begin{figure*}[t]
  \centering
  % TODO: ใส่รูป system architecture diagram
  \fbox{\parbox{0.95\textwidth}{
    \centering
    \vspace{2cm}
    \textbf{[Figure 1: TempoRAG-KG System Architecture]}\\[0.5em]
    \textit{Offline pipeline (top):} Documents $\rightarrow$ Chunking $\rightarrow$
    LLM Triple Extraction + temporal tagging $\rightarrow$ Neo4j KG + FAISS Index\\[0.3em]
    \textit{Online pipeline (bottom):} Query $\rightarrow$ Temporal year detection
    $\rightarrow$ Seed retrieval $\rightarrow$ \textbf{Temporal filter}
    $\rightarrow$ GEAR beam search $\rightarrow$ \textbf{GoG fill-in}
    $\rightarrow$ MST context organization $\rightarrow$ LLM $\rightarrow$ Answer\\[0.3em]
    \textit{Highlight boxes:} TempoRAG-KG additions shown in \textcolor{blue}{blue};
    KG$^2$RAG components shown in gray
    \vspace{2cm}
  }}
  \caption{Overview of the TempoRAG-KG framework. Components highlighted in
  blue are our additions to the KG$^2$RAG backbone. The offline pipeline
  constructs a temporally-annotated KG; the online pipeline applies temporal
  filtering before graph expansion and GoG fill-in for incomplete facts.}
  \label{fig:architecture}
\end{figure*}

\subsection{Datasets}

\paragraph{HotpotQA.}
We use the distractor development set of HotpotQA \cite{yang2018hotpotqa},
comprising 7,405 two-hop questions drawn from Wikipedia with ten candidate
paragraphs each. Our EDA reveals that \textbf{35.0\%} of questions
(2,593/7,405) contain temporal expressions. Bridge-type questions exhibit
higher temporal prevalence (39\%) than comparison-type questions (21\%),
consistent with the expectation that multi-hop reasoning more often
requires temporal grounding. Table~\ref{tab:eda} summarizes the
distribution by temporal pattern type.

\paragraph{MuSiQue.}
MuSiQue \cite{trivedi2022musique} provides 2,417 development questions
with 2--4 reasoning hops, designed to preclude shortcut reasoning.
\textbf{38.3\%} of questions (926/2,417) contain temporal expressions.
Crucially, temporal prevalence increases monotonically with hop count:
30.2\% for 2-hop, 47.2\% for 3-hop, and 46.7\% for 4-hop questions,
empirically motivating temporal-aware retrieval for complex reasoning chains.

\paragraph{Temporal Extraction Gold Standard.}
For RQ3, we manually annotate 200 passages sampled from the HotpotQA
training set, recording ground-truth $[\mathtt{valid\_from},
\mathtt{valid\_to}]$ for all extractable triples. Passages are stratified
across three temporal expression types: explicit year mention (e.g.,
``became CEO in 2011''), implicit/relative (e.g., ``during his tenure''),
and conflicting (multiple dates for the same relation).

% EDA Table
\begin{table}[t]
  \centering
  \small
  \begin{tabular}{lrr}
    \toprule
    \textbf{Temporal Pattern} & \textbf{HotpotQA} & \textbf{MuSiQue} \\
    \midrule
    Year mention        & 17.8\% & 6.7\%  \\
    When/what year      & 7.3\%  & 20.6\% \\
    Before/after        & 4.2\%  & 7.1\%  \\
    First/last/earliest & 7.7\%  & 7.1\%  \\
    Time period         & 2.3\%  & 3.9\%  \\
    Age/duration        & 0.7\%  & 1.8\%  \\
    \midrule
    \textbf{Any temporal} & \textbf{35.0\%} & \textbf{38.3\%} \\
    \bottomrule
  \end{tabular}
  \caption{Temporal expression prevalence in HotpotQA and MuSiQue
  development sets. Questions may match multiple patterns.}
  \label{tab:eda}
\end{table}

\subsection{Data Preprocessing}

\paragraph{Chunking.}
Each document is segmented into overlapping chunks of 512 tokens with a
100-token overlap, following KG$^2$RAG \cite{zhu2025kg2rag}. Each chunk
retains a reference to its source document.

\paragraph{Triple Extraction with Temporal Annotation.}
We prompt GPT-4 to extract structured triples from each chunk in the
following format: $(\mathtt{head}, \mathtt{relation}, \mathtt{tail},
\mathtt{valid\_from}, \mathtt{valid\_to})$. The prompt instructs the model
to infer temporal validity from contextual cues (explicit years, relative
temporal expressions, succession language). If no temporal information is
available, fields are set to \texttt{null}. Figure~\ref{fig:architecture}
illustrates this extraction step.

\paragraph{KG Storage and Indexing.}
Extracted triples are stored in Neo4j with $\mathtt{valid\_from}$ and
$\mathtt{valid\_to}$ as edge properties. Entity names are embedded using
\texttt{all-MiniLM-L6-v2} \cite{reimers2019sentencebert} and stored in a
FAISS index for efficient similarity search.

\subsection{TempoRAG-KG Pipeline}

The online pipeline proceeds in six steps:

\textbf{(1) Temporal Year Detection.}
We apply a regex-based classifier to identify whether the input query
contains a temporal expression and, if so, extract the target year
$y_q$.

\textbf{(2) Seed Retrieval.}
The query is embedded and the top-5 most similar chunks are retrieved
from the FAISS chunk index, following KG$^2$RAG.

\textbf{(3) Temporal Filtering.}
For each candidate edge $e$ with properties $v_f$ ($\mathtt{valid\_from}$)
and $v_t$ ($\mathtt{valid\_to}$), we retain $e$ if and only if:
\begin{equation}
  (v_f = \texttt{null} \;\lor\; v_f \leq y_q) \;\land\;
  (v_t = \texttt{null} \;\lor\; v_t \geq y_q)
  \label{eq:filter}
\end{equation}
Edges without temporal metadata are retained conservatively.

\textbf{(4) GEAR Beam Search.}
Graph expansion proceeds via beam search \cite{ma2025gear} with beam
size $b=3$ and maximum depth $d=2$. At each step, candidate neighbors
are scored by cosine similarity to the current subgoal embedding, avoiding
LLM calls during traversal.

\textbf{(5) GoG Fill-in.}
When expansion reaches a dead end (no outgoing edges satisfy the temporal
filter), we invoke a GoG-style LLM fill-in \cite{he2024gog}: the model is
prompted to supply the missing fact from parametric knowledge, along with
a confidence score. Fill-in facts with confidence below 0.7 are discarded.

\textbf{(6) MST Context Organization and Generation.}
Retrieved subgraph facts are organized via Maximum Spanning Tree pruning
(weighted by semantic relevance) into coherent paragraphs, which are
prepended to the generation prompt for LLaMA-3-8B or GPT-4.

% -------------------------------------------------------
% Main Results Table
% -------------------------------------------------------
\begin{table*}[t]
  \centering
  \small
  \begin{tabular}{lcccccc}
    \toprule
    \multirow{2}{*}{\textbf{Model}} &
    \multicolumn{2}{c}{\textbf{Overall}} &
    \multicolumn{2}{c}{\textbf{Temporal Subset}} &
    \multicolumn{2}{c}{\textbf{Non-temporal Subset}} \\
    \cmidrule(lr){2-3}\cmidrule(lr){4-5}\cmidrule(lr){6-7}
    & \textbf{F1} & \textbf{EM} & \textbf{F1} & \textbf{EM} & \textbf{F1} & \textbf{EM} \\
    \midrule
    Vanilla RAG         & -- & -- & -- & -- & -- & -- \\
    GraphRAG            & -- & -- & -- & -- & -- & -- \\
    KG$^2$RAG           & -- & -- & -- & -- & -- & -- \\
    \midrule
    \textbf{TempoRAG-KG (ours)} & \textbf{--} & \textbf{--} & \textbf{--} & \textbf{--} & \textbf{--} & \textbf{--} \\
    \bottomrule
  \end{tabular}
  \caption{Main results on HotpotQA distractor dev set. Results reported
  as F1 and Exact Match (EM). Temporal and non-temporal subsets contain
  2,593 and 4,812 questions respectively. \textbf{Bold} = best result.
  (Results to be filled after experiments.)}
  \label{tab:main_results}
\end{table*}

% -------------------------------------------------------
% Ablation Scores Table
% -------------------------------------------------------
\begin{table}[t]
  \centering
  \small
  \begin{tabular}{lcc}
    \toprule
    \textbf{Ablation Condition} &
    \textbf{Overall F1} &
    \textbf{Temporal F1} \\
    \midrule
    Vanilla RAG                    & -- & -- \\
    KG$^2$RAG (backbone)           & -- & -- \\
    \quad +Temporal filter only    & -- & -- \\
    \quad +GoG fill-in only        & -- & -- \\
    \quad \textbf{Full model (ours)} & \textbf{--} & \textbf{--} \\
    \midrule
    $\Delta$ vs.\ KG$^2$RAG       & --  & -- \\
    \bottomrule
  \end{tabular}
  \caption{Ablation study on HotpotQA. Each row removes one component
  from the full model. $\Delta$ = improvement over KG$^2$RAG backbone.
  (Results to be filled after experiments.)}
  \label{tab:ablation_scores}
\end{table}

% -------------------------------------------------------
% Temporal Extraction Table (RQ3)
% -------------------------------------------------------
\begin{table}[t]
  \centering
  \small
  \begin{tabular}{lccc}
    \toprule
    \textbf{Pattern Type} & \textbf{Prec.} & \textbf{Recall} & \textbf{F1} \\
    \midrule
    Explicit year mention   & -- & -- & -- \\
    Implicit / relative     & -- & -- & -- \\
    Conflicting dates       & -- & -- & -- \\
    \midrule
    \textbf{Overall}        & \textbf{--} & \textbf{--} & \textbf{--} \\
    \bottomrule
  \end{tabular}
  \caption{Temporal extraction accuracy on 200 manually annotated
  HotpotQA passages (RQ3). Precision and recall computed against
  gold $[\mathtt{valid\_from}, \mathtt{valid\_to}]$ annotations,
  with $\pm$1 year tolerance for boundary matching.
  (Results to be filled after annotation.)}
  \label{tab:extraction}
\end{table}

\subsection{Baselines}

We compare against four baselines:
\begin{itemize}
  \item \textbf{Vanilla RAG}: Dense retrieval with GPT-4 generation; no graph.
  \item \textbf{GraphRAG}: Microsoft community-summary graph RAG \cite{edge2024graphrag}.
  \item \textbf{KG$^2$RAG}: Reproduced NAACL 2025 baseline \cite{zhu2025kg2rag}.
  \item \textbf{TempoRAG-KG (ours)}: Full system with temporal filter and GoG fill-in.
\end{itemize}

\subsection{Ablation Design}

To isolate each component's contribution (RQ2), we evaluate five conditions:

\begin{table}[t]
  \centering
  \small
  \begin{tabular}{lccc}
    \toprule
    \textbf{Condition} & \textbf{Temporal} & \textbf{GoG} & \textbf{Graph Exp.} \\
    \midrule
    Vanilla RAG            & \xmark & \xmark & \xmark \\
    KG$^2$RAG              & \xmark & \xmark & \cmark \\
    +Temporal only         & \cmark & \xmark & \cmark \\
    +GoG only              & \xmark & \cmark & \cmark \\
    \textbf{Full (ours)}   & \cmark & \cmark & \cmark \\
    \bottomrule
  \end{tabular}
  \caption{Ablation conditions. \cmark = component enabled;
  \xmark = component disabled.}
  \label{tab:ablation}
\end{table}

\subsection{Evaluation Metrics}

We report:
\begin{itemize}
  \item \textbf{F1 and EM}: Token-level F1 and Exact Match, computed
    overall and separately for temporal and non-temporal question subsets.
  \item \textbf{Retrieval Precision@5}: Fraction of retrieved chunks
    containing the gold supporting facts.
  \item \textbf{Temporal Extraction F1}: Precision and recall of
    $[\mathtt{valid\_from}, \mathtt{valid\_to}]$ against gold annotation
    (RQ3 only; tolerates $\pm$1 year for boundary matching).
  \item \textbf{Inference Time}: Average seconds per query.
\end{itemize}

% -------------------------------------------------------
% 5. EXPECTED RESULTS
% -------------------------------------------------------
\section{Expected Results and Discussion}

\subsection{RQ1: Temporal Failure Mode Distribution}

We expect stale-fact retrieval (Type 1) to constitute the majority of
temporal errors in KG$^2$RAG, as the most common temporal pattern in
HotpotQA is explicit year mention (17.8\%), which frequently surfaces
conflicting facts from different eras. Type 4 (temporal hop failure) is
expected to be rare in HotpotQA (2-hop only) but more frequent in
MuSiQue 4-hop questions. This taxonomy, to our knowledge, is the first
systematic categorization of temporal failures in general-domain KG-RAG.

\subsection{RQ2: Impact of Temporal Tagging}

We hypothesize that the full TempoRAG-KG model will outperform
KG$^2$RAG on the temporal question subset by a meaningful margin,
while maintaining comparable performance on non-temporal questions
(confirming that temporal filtering does not harm non-temporal retrieval).
We further expect the ablation results to confirm that temporal filtering
and GoG fill-in contribute independently and additively, as they address
distinct failure modes (Types 1--3 vs.\ incomplete KG facts respectively).

\begin{figure}[t]
  \centering
  \fbox{\parbox{0.45\textwidth}{
    \centering
    \vspace{1.5cm}
    \textbf{[Figure 3: Expected Ablation Results]}\\[0.5em]
    Bar chart comparing F1 scores (overall + temporal subset):\\
    Vanilla RAG $<$ KG$^2$RAG $<$ +Temporal $<$ Full Model\\[0.3em]
    Secondary plot: non-temporal F1 stays flat\\
    across all conditions (temporal filter has no\\
    negative effect on non-temporal questions)
    \vspace{1.5cm}
  }}
  \caption{Expected ablation results on HotpotQA. We predict gains
  concentrated in the temporal question subset, with non-temporal F1
  remaining stable.}
  \label{fig:ablation}
\end{figure}

\subsection{RQ3: Temporal Extraction Accuracy}

We expect high precision on explicit year-mention expressions ($>$80\%)
but substantially lower recall on implicit/relative expressions (e.g.,
``during his tenure'', ``at the time''), where the model must infer
temporal boundaries from context. Conflicting date expressions (where
a text contains multiple dates for the same entity-relation pair) are
expected to pose the greatest challenge. This characterization directly
informs the practical ceiling of TempoRAG-KG's temporal filtering
quality.

% -------------------------------------------------------
% 6. PROGRESS SUMMARY
% -------------------------------------------------------
\section{Progress Summary}

\paragraph{Completed.}
\begin{itemize}
  \item Literature review: 16 papers across four team members
    (4 papers each, following course requirements).
  \item EDA of HotpotQA (7,405 questions: 35.0\% temporal) and
    MuSiQue (2,417 questions: 38.3\% temporal) using automated
    regex-based temporal classifier with seven pattern categories.
  \item Temporal failure taxonomy: four types proposed based on
    qualitative analysis of KG-RAG failure modes.
  \item Identification and analysis of closest competitor (EvoReasoner;
    arXiv 2509.15464) and clear differentiation on three dimensions:
    explicit vs.\ implicit temporal model, general-domain vs.\
    structured KG, and benchmark choice.
\end{itemize}

\paragraph{In Progress.}
\begin{itemize}
  \item Neo4j schema design and setup with temporal edge properties.
  \item Temporal extraction prompt engineering and pilot evaluation
    on 20 HotpotQA passages.
\end{itemize}

\paragraph{Remaining Work.}
\begin{itemize}
  \item Reproduction of KG$^2$RAG baseline on HotpotQA distractor dev set.
  \item Gold standard annotation of 200 passages for RQ3.
  \item Full pipeline implementation: temporal filter, GEAR beam search,
    GoG fill-in integration.
  \item Ablation experiments (Table~\ref{tab:ablation}).
  \item Streamlit demo with edge case scenarios.
\end{itemize}

\paragraph{Potential Concerns.}
The primary risk is temporal extraction quality on implicit expressions.
If extraction F1 on implicit patterns falls below 60\%, the temporal
filter may introduce false negatives (discarding valid facts). We have
a mitigation strategy: conservatively retaining triples with
$\mathtt{valid\_from} = \mathtt{null}$ (Equation~\ref{eq:filter}),
ensuring that extraction errors degrade gracefully rather than
catastrophically.

% -------------------------------------------------------
% 7. CONCLUSION
% -------------------------------------------------------
\section{Conclusion}

We have presented TempoRAG-KG, a temporally-aware Knowledge Graph-augmented
RAG framework that addresses a systematic and underexplored failure mode in
existing KG-RAG systems: the inability to distinguish facts valid at
different points in time. By attaching explicit validity intervals to KG
triples and filtering graph expansion temporally, TempoRAG-KG targets the
35\% of HotpotQA and 38\% of MuSiQue questions that require temporal
reasoning. Our temporal failure taxonomy and extraction study provide
the diagnostic groundwork needed to understand when and why temporal
reasoning fails in graph-augmented systems. This work contributes to the
growing body of research at the intersection of structured knowledge
retrieval and temporal reasoning, and opens a tractable path toward
knowledge systems that remain reliably accurate across time.

% -------------------------------------------------------
% REFERENCES
% -------------------------------------------------------
\bibliography{references}

\appendix

% -------------------------------------------------------
% APPENDIX A — Temporal Extraction Prompt
% -------------------------------------------------------
\section{Temporal Extraction Prompt}
\label{app:prompt}

\begin{verbatim}
You are a knowledge extraction assistant.
Given the following text passage, extract all
factual triples. For each triple, also infer
temporal validity if mentioned or implied.

Format each triple as JSON:
{
  "head": "entity name",
  "relation": "relation type",
  "tail": "entity or value",
  "valid_from": year (int) or null,
  "valid_to": year (int) or null,
  "confidence": float between 0 and 1
}

Rules:
- Only extract explicitly stated facts
- For valid_from/valid_to: use null if unknown
- If a fact is described as ongoing, set
  valid_to to null
- If two dates conflict for the same triple,
  extract both as separate triples

Text: {chunk_text}
\end{verbatim}}

% -------------------------------------------------------
% APPENDIX B — EDA Details
% -------------------------------------------------------
\section{EDA Methodology}
\label{app:eda}

Temporal questions were identified using seven regex patterns applied
to the lowercased question string. A question was classified as temporal
if it matched at least one pattern. The patterns target: (1) explicit
year mentions (\texttt{\textbackslash b(1[5-9]\textbackslash d\{2\}|20\textbackslash d\{2\})\textbackslash b}),
(2) when/what-year phrases, (3) before/after/since, (4) time period
expressions, (5) first/last/earliest superlatives, (6) age and duration
phrases, and (7) temporal relation words (predecessor, successor, etc.).

\end{document}
